% Irodov, Arihant PDF pp. 26--27, printed pp. 17--18.
\exproblem{101}{17}
子弹穿过厚度为 $h$ 的木板，速率由 $v_0$ 降为 $v$。
若阻力大小与速率平方成正比，求子弹在板内运动的时间。
\exsolution
令阻力为 $F=-kv^2$，则
$m\dot v=-kv^2$。利用
$\dot v=v\,\mathrm dv/\mathrm dx$ 可得
\[
m v\frac{\mathrm dv}{\mathrm dx}=-kv^2
\ \Longrightarrow\
\frac{\mathrm dv}{v}=-\frac km\,\mathrm dx.
\]
从入板到出板积分，
$\ln(v_0/v)=kh/m$。另一方面
\[
\frac{\mathrm d}{\mathrm dt}\left(\frac1v\right)=\frac km,
\quad
t=\frac mk\left(\frac1v-\frac1{v_0}\right).
\]
消去未知 $k/m$：
\[
\boxed{t=\frac{h(v_0-v)}
 {v_0v\ln(v_0/v)}}.
\]
\excheck 当 $v\to v_0$，对数展开使 $t\to h/v_0$，
与几乎无阻力时的通过时间一致。

\exproblem{102}{17}
物块从静止开始沿倾角 $\alpha$ 的斜面下滑，
摩擦系数随已走过的距离 $x$ 增长为 $k(x)=ax$，$a>0$。
求物块直到停止的总路程，以及全过程中的最大速率。
\exsolution
沿斜面方向列方程
$m\dot v=mg(\sin\alpha-ax\cos\alpha)$。
取 $\dot v=v\,\mathrm dv/\mathrm dx$ 并积分初值 $v(0)=0$：
\[
\frac{v(x)^2}{2}
=g\left(x\sin\alpha-\frac12ax^2\cos\alpha\right).
\]
除 $x=0$ 之外，$v=0$ 时的路程为
\[
\boxed{s_{\rm stop}=\frac{2\tan\alpha}{a}}.
\]
速率最大时 $\dot v=0$，即
$x_m=\tan\alpha/a$，代回得
\[
\boxed{v_{\max}=
\sqrt{\frac ga\sin\alpha\tan\alpha}}.
\]
\excheck 在 $x=x_m$ 处加速度由正变负；
在 $s_{\rm stop}$ 前 $v^2\ge0$，能量表达式始终自洽。

\exproblem{103}{17}
质量为 $m$ 的物块静置于水平粗糙面上，摩擦系数为 $k$。
从 $t=0$ 开始施加随时间线性增长的水平力 $F=at$，$a>0$。
求前 $t$ 秒走过的距离。
\exsolution
力尚未超过最大静摩擦力 $kmg$ 时物块静止，
临界时刻
\[
t_0=\frac{kmg}{a}.
\]
当 $t>t_0$，按静、动摩擦系数相同的原书模型，
净水平力为 $a(t-t_0)$，故
\[
m\ddot x=a(t-t_0),\quad
v(t)=\frac{a}{2m}(t-t_0)^2,\quad
s(t)=\frac{a}{6m}(t-t_0)^3.
\]
因此
\[
\boxed{s(t)=
\begin{cases}
0,&0\le t\le t_0,\\[2pt]
\dfrac{a(t-t_0)^3}{6m},&t>t_0.
\end{cases}}
\]
\excheck 在 $t=t_0$ 处，位移、速度、加速度均连续。
若静摩擦系数与动摩擦系数不同，应分别引入两者，
本表达式不再自动适用。

\exproblem{104}{17}
质量 $m$ 的物体由地面以初速 $v_0$ 竖直上抛，
空气阻力总与速度反向，大小为 $kv^2$，$k>0$。
问物体回到原抛出高度时向下的速率 $v'$。
\exsolution
上升阶段取向上位移 $y$、速率 $u$，
重力和阻力均向下：
\[
m u\frac{\mathrm du}{\mathrm dy}=-mg-ku^2.
\]
从 $u=v_0$ 积分到顶点 $u=0$，最高高度为
\[
H=\int_0^{v_0}\frac{mu\,\mathrm du}{mg+ku^2}
=\frac{m}{2k}\ln\left(1+\frac{kv_0^2}{mg}\right).
\]
下落阶段取向下路程 $z$、速率 $u$，
运动方程 $m u\,\mathrm du/\mathrm dz=mg-ku^2$。
从 $u(0)=0$ 得
\[
u(z)^2=\frac{mg}{k}\left(1-e^{-2kz/m}\right).
\]
令 $z=H$，得
\[
\boxed{v'=\frac{v_0}{\sqrt{1+kv_0^2/(mg)}}}.
\]
\excheck 该速率严格小于 $v_0$，且 $k\to0$ 时趋于 $v_0$，
符合空气阻力耗散机械能的事实。

\exproblem{105}{17}
质量为 $m$ 的质点在平面内受大小恒为 $F$ 的力，
力矢量以角速度 $\omega>0$ 匀速旋转。
质点在 $t=0$ 时静止。求：
(a) 速度矢量随时间变化；
(b) 相邻两次停止间的路程及平均速率。
\exsolution
选坐标轴使初始外力沿 $x$ 正向，
$\bm F(t)=F(\cos\omega t\,\bm i+
\sin\omega t\,\bm j)$。由 $\bm v(0)=0$ 积分：
\[
\boxed{\bm v(t)=\frac{F}{m\omega}
[\sin\omega t\,\bm i+(1-\cos\omega t)\bm j]}.
\]
速率为
\[
v(t)=\frac{2F}{m\omega}
\left|\sin\frac{\omega t}{2}\right|.
\]
相邻停止时刻 $t=0$ 与 $t=2\pi/\omega$，
故
\[
\boxed{\Delta s=\int_0^{2\pi/\omega}v(t)\,\mathrm dt
=\frac{8F}{m\omega^2},\qquad
\overline{v}_{\rm speed}=\frac{\Delta s}{2\pi/\omega}
=\frac{4F}{\pi m\omega}}.
\]
\excheck 此处平均\emph{速率}是书末答案中的 $4F/(\pi m\omega)$；
若求平均\emph{速度矢量}，则需位移除以时间，
其值为 $(F/m\omega)\bm j$，二者不相同。

\exproblem{106}{18}
图1.27中，一小圆片在倾角 $\alpha$ 的斜面上滑动，
摩擦系数 $k=\tan\alpha$。
其速度与斜面最大下降方向 $x$ 的夹角为 $\varphi$，
初始 $\varphi_0=\pi/2$，初速为 $v_0$。
求速率 $v$ 关于 $\varphi$ 的表达式。
\exsolution
沿斜面最大下降方向取 $\bm i$，
重力沿斜面分量大小 $g\sin\alpha$，
摩擦加速度大小 $kg\cos\alpha=g\sin\alpha$，
且始终与速度相反。
设 $b=g\sin\alpha$，速度方向单位矢量
$(\cos\varphi,\sin\varphi)$。动力学沿速度切线和法线分解：
\[
\dot v=b(\cos\varphi-1),\qquad
v\dot\varphi=-b\sin\varphi.
\]
消去时间：
\[
\frac{\mathrm dv}{v}
=\frac{1-\cos\varphi}{\sin\varphi}\,\mathrm d\varphi
=\tan\frac\varphi2\,\mathrm d\varphi.
\]
积分并用 $v(\pi/2)=v_0$：
\[
\ln v=-2\ln\cos(\varphi/2)+\ln(v_0/2),
\quad
\boxed{v(\varphi)=\frac{v_0}{1+\cos\varphi}}.
\]
\excheck 初始 $\varphi=\pi/2$ 时恢复 $v_0$；
后续速度逐渐转向下坡方向 $\varphi\to0$，
极限速率为 $v_0/2$。

\exproblem{107}{18}
一条总长 $\ell$、质量均匀的链条铺在光滑球面上，
球半径 $R$，上端固定在顶点，且 $\ell<\pi R/2$。
解除上端约束的瞬间，求链条各元的切向加速度。
\exsolution
设链条从球顶沿同一子午线向下延伸，
取弧长坐标 $s\in[0,\ell]$，球顶对应 $s=0$。
绳不可伸长，各元沿链条的瞬时切向加速度相同，记为 $a$。
沿整条链对切向方向求受力积分，内部拉力彼此抵消，
光滑球面的支持力无切向分量。线密度为
$\lambda=m/\ell$，弧长 $s$ 处重力切向分量为
$\lambda g\sin(s/R)\,\mathrm ds$。于是
\[
ma=\lambda g\int_0^\ell\sin(s/R)\,\mathrm ds
=\lambda gR[1-\cos(\ell/R)].
\]
故
\[
\boxed{a=\frac{gR}{\ell}[1-\cos(\ell/R)]},
\]
所有链元沿子午线向下同加速度。
\excheck $\ell\ll R$ 时
$a\simeq g\ell/(2R)$；链很短时仅有很小的平均切向重力。

\exproblem{108}{18}
一小物块放在半径 $R$ 的光滑球体顶端，
随后球体突然开始以恒定水平加速度 $a_0$ 平移，
物块相对球体滑下。求：
(a) 刚离面时相对球体的速率；
(b) 离面位置半径与竖直向上方向的夹角 $\theta_0$；
并计算 $a_0=g$ 时的角度。
\exsolution
在随球体加速的非惯性系中，物块受重力和恒定水平
惯性力 $m a_0$，方向与球体加速方向相反。
记 $\theta$ 为向\emph{惯性力一侧}的偏角，
坐标为 $x=R\sin\theta$（向惯性力一侧）、$y=R\cos\theta$。
从顶部开始相对静止，保守有效力场的能量式为
\[
\frac{v'^2}{2}=gR(1-\cos\theta)+a_0R\sin\theta.
\]
向球心投影牛顿第二定律：
\[
\frac{v'^2}{R}=g\cos\theta-a_0\sin\theta-\frac Nm.
\]
离面时 $N=0$。联立消去 $v'$：
\[
g\cos\theta_0-a_0\sin\theta_0=\frac23g.
\]
因此
\[
\boxed{v'_{\rm detach}=\sqrt{\frac{2gR}{3}}}
\]
并设 $\eta=a_0/g$，解
$\cos\theta_0-\eta\sin\theta_0=2/3$ 得
\[
\boxed{\cos\theta_0=
\frac{2+\eta\sqrt{5+9\eta^2}}{3(1+\eta^2)}}.
\]
当 $\eta=1$ 时
$\cos\theta_0=(2+\sqrt{14})/6$，
故 $\boxed{\theta_0\simeq16.9^\circ}$。
\excheck $\eta\to0$ 得 $\cos\theta_0=2/3$，
退化为静止球面上的第87题。这里的速率是\emph{相对球}速率。

\exproblem{109}{18}
质点在平面内受一始终垂直于其速度的力，
其大小与到某固定点距离 $r$ 的 $n$ 次方成反比，
$|\bm F|=C/r^n$。求哪些 $n$ 使圆轨道运动稳定。
\exsolution
因 $\bm F\perp\bm v$，速率 $v$ 恒定。
以固定点为极点，圆轨道半径 $R$ 处
$C/R^n=mv^2/R$，即 $C=mv^2R^{n-1}$。
考虑小的径向扰动 $r=R+\delta r$。
垂直速度且指向轨迹凹侧的力，其极坐标分量在线性近似为
\[
F_r\simeq-\frac{C}{r^n},\qquad
F_\varphi\simeq\frac{C}{r^n}\frac{\dot r}{v}.
\]
对角动量应用
$m\,\mathrm d(r^2\dot\varphi)/\mathrm dt=rF_\varphi$，
积分到一阶得到
$\delta(r^2\dot\varphi)=v\,\delta r$，
所以 $\dot\varphi\simeq(v/R)(1-\delta r/R)$。
将这些代入径向运动方程
$m(\ddot r-r\dot\varphi^2)=F_r$，保留一阶项：
\[
\delta\ddot r+\frac{(1-n)v^2}{R^2}\,\delta r=0.
\]
因而稳定的小振动要求恢复系数为正，即
\[
\boxed{n<1}.
\]
\excheck $n=1$ 为线性中性边界，$n>1$ 径向扰动指数增长。
推导已使用“力始终与速度垂直”而非中央力假设，
不能将开普勒中央力稳定性判据直接挪用。

\exproblem{110}{18}
图1.28所示半径为 $R$ 的光滑半圆形竖直圆弧杆
绕过圆心的竖直轴 $OO'$ 以角速度 $\omega$ 恒速旋转。
质量为 $m$ 的小环可在圆弧上自由滑动。
求稳态平衡位置对应的角度 $\theta$，并讨论稳定性。
\exsolution
从圆弧最低点起量角 $\theta$，小环距转轴
$r=R\sin\theta$，高度取 $y=-R\cos\theta$。
在旋转参考系中，沿圆弧的有效势能为
\[
U_{\rm eff}(\theta)=
-mgR\cos\theta-\frac12m\omega^2R^2\sin^2\theta.
\]
平衡位置满足
\[
\frac{\mathrm dU_{\rm eff}}{\mathrm d\theta}
=mR\sin\theta(g-\omega^2R\cos\theta)=0.
\]
故始终有最低点
$\boxed{\theta_1=0}$；
仅当 $\omega^2R\ge g$ 时还有
\[
\boxed{\theta_2=\arccos\frac{g}{\omega^2R}}.
\]
检查有效势能二阶导数：
当 $\omega^2R<g$ 时 $\theta_1$ 稳定；
当 $\omega^2R>g$ 时 $\theta_1$ 失稳而
$\theta_2$ 稳定；临界 $\omega^2R=g$ 处需要考虑高阶项。
\excheck 这是旋转参考系的稳态，而不是小环相对地面静止；
平衡位置随 $\omega^2$ 改变，旋转方向无关。
